% CLAUDE REPORT::WORSTCASE-BRIDGE-0001::2026-11-11T12:48-08:00
% !TEX program = xelatex
% Compiles with XeLaTeX or pdfLaTeX (ASCII source).
\documentclass[11pt]{article}
\usepackage[margin=1in]{geometry}
\usepackage{amsmath,amssymb,amsthm}
\usepackage{booktabs}
\usepackage{xurl}
\newtheorem{theorem}{Theorem}
\newtheorem{lemma}[theorem]{Lemma}
\newtheorem{proposition}[theorem]{Proposition}
\newtheorem{corollary}[theorem]{Corollary}
\theoremstyle{definition}
\newtheorem{definition}[theorem]{Definition}
\theoremstyle{remark}
\newtheorem*{remark}{Remark}
\newcommand{\SAT}{\mathrm{SAT}}
\newcommand{\UNSAT}{\mathrm{UNSAT}}
\newcommand{\poly}{\operatorname{poly}}
\newcommand{\cost}{\operatorname{cost}}
\newcommand{\Wb}{\overline W}
\newcommand{\status}[1]{\textbf{[#1]}}

\title{CLAUDE WORSTCASE-BRIDGE-0001\\[2pt]
\large The worst-case derivation target $\Wb_R(N)$ and the bridge
$\SAT\in\mathrm P\Rightarrow\Wb_R(N)\le\poly(N)$:\\
a two-sided squeeze, a reflection-system bridge, and five-role assignments}
\author{CLAUDE (scheduled research routine; not an LSPV role)}
\date{2026-11-11 12:48 America/Los\_Angeles}

\begin{document}
\maketitle

\noindent\textbf{Status.} \textbf{P versus NP remains OPEN.} Every load-bearing
arrow toward $\mathrm P\ne\mathrm{NP}$ in this report is OPEN. The lemmas
labelled PROVED/EXACT are elementary and hold in the stated models; most of
them are folklore in propositional proof complexity, and we claim no
originality for them, only that they set the terms of the new shared target.
This is an additive report. It changes no LSPV schedule, role, ownership
region or allocation, and it does not edit root \texttt{main.tex}; LSPV-5
decides whether to absorb it.

\section{The target, stated exactly}

Fix a sound and complete refutation calculus $R$ for CNF that is a
Cook--Reckhow proof system: given $(F,\pi)$, a polynomial-time verifier checks
that $\pi$ is an $R$-refutation of $F$. The user's multisign DAG hypergraph
reduction system is an instance: states are signed hyperedge sets, a
\emph{reduction step} is a sound local rule, and a \emph{reduction sequence}
that reaches the empty (contradictory) state is a refutation. Let
$\cost_R(\pi)$ be the cost measure: DAG nodes counted once, plus encoding
bits, and optionally peak memory. Define
\[
 w_R(F)=\min_{\pi:\,F\vdash_R\bot}\cost_R(\pi),\qquad
 \Wb_R(N)=\max_{\substack{F\in\UNSAT\\|\langle F\rangle|\le N}} w_R(F).
\]
This is the same object as $W_R$ in the Codex audit
\texttt{research/agents/codex/2026-11-11\_1203\_SAT\_POLYNOMIAL\_ASSUMPTION\_AUDIT.tex}.
Two properties are exact:

\begin{itemize}
\item \emph{Free choice of sequence is sound for lower bounds.} \status{PROVED/EXACT}
$w_R(F)$ is a minimum over all valid sequences, so any procedure that
\emph{executes} an $R$-refutation of $F$ pays at least $w_R(F)$, however the
sequence was found. Search can only add cost. This confirms the user's
clarification.
\item \emph{The minimum is a nondeterministic measure.} \status{PROVED/EXACT}
Because the optimal sequence is free, $\Wb_R(N)\le\poly(N)$ says exactly that
$R$ is \emph{polynomially bounded} (when $\cost_R$ is polynomially related to
proof length). By Cook--Reckhow, a polynomially bounded system exists if and only
if $\mathrm{NP}=\mathrm{coNP}$ \cite{cookreckhow1979}.
\end{itemize}

So a superpolynomial lower bound on $\Wb_R$ is a statement in proof complexity.
It reaches the time of an \emph{arbitrary} SAT algorithm only through the
bridge
\[
 \mathrm{Bridge}(R):\qquad \SAT\in\mathrm P\ \Longrightarrow\ \Wb_R(N)\le\poly(N),
\]
together with
\[
 \mathrm{LB}(R):\qquad \Wb_R(N)\ne N^{O(1)}.
\]
$\mathrm{Bridge}(R)\wedge\mathrm{LB}(R)\Rightarrow\SAT\notin\mathrm P$
\status{PROVED/EXACT, by contraposition}. Both conjuncts are OPEN for every
$R$ for which neither is already known to be equivalent to P versus NP itself;
see Section~\ref{sec:squeeze}.

\section{CORPUS REVIEWED}
Repository tree survey; \texttt{AGENTS.md} (including the 2026-11-09
knot/link handoff and the 2026-10-24 all-ARC/FLMV allocation); \texttt{CODEX.md};
root \texttt{main.tex} (heading and status lines only, since this run does not
edit it); the Codex 2026-11-11 12:03 SAT-polynomial-assumption audit (its
definition of $W_R$, the global bridge (B), the warning that
``Cook--Reckhow is existential over systems'', and the open-status table);
the newest LSPV coordination handoffs
(LSPV-1/573, LSPV-2/558, LSPV-3/570--571, LSPV-4/543--544, LSPV-5/488R1),
whose subjects are native calculus/CLI replay and resource accounting for fixed
W61 proofs; the earlier Claude ledger \texttt{research/agents/claude/world-ledger/}
(row L13, Buss--Williams); and repository search for ``worst-case directive'',
``$\overline W_R$'' and ``weakest sufficient bridge''.

\section{DUPLICATION / CONTRADICTION AUDIT}
\begin{itemize}
\item Repository search found \emph{no} existing file for the ``shared worst-case
research directive'' that the scheduled prompt refers to. This run creates it as
\texttt{research/coordination/2026-11-11\_1248\_WORST\_CASE\_DERIVATION\_DIRECTIVE\_LIVING\_LOG.md}.
\item The Codex 12:03 audit already defines $W_R$ and studies a
quasipolynomial, all-decider version of the bridge. This report studies the
\emph{polynomial} bridge, adds the monotone squeeze (Lemma~\ref{lem:mono}), the
reformulation trap (Proposition~\ref{prop:trap}) and the reflection system
(Theorem~\ref{thm:refl}). It does not contradict the audit. In particular it
agrees that Cook--Reckhow by itself gives no bridge for a \emph{chosen} $R$.
\item The LSPV W61 and native-calculus work certifies \emph{upper} bounds, that
is, explicit short derivations and their resource cost, for specific
instances. Such certificates are useful for the Red Team, because they bound
$w_R$ from above, but they say nothing about $\mathrm{LB}(R)$.
\end{itemize}

\section{The two-sided squeeze}\label{sec:squeeze}

\begin{lemma}[Monotonicity under simulation]\label{lem:mono}
\status{PROVED/EXACT} Suppose $R'$ $p$-simulates $R$, meaning that every
$R$-refutation $\pi$ of $F$ maps to an $R'$-refutation of cost
$\le q(|F|+\cost_R(\pi))$ for a fixed polynomial $q$. Then
$\Wb_{R'}(N)\le q(N+\Wb_R(N))$. Consequently
$\mathrm{Bridge}(R)\Rightarrow\mathrm{Bridge}(R')$ and
$\mathrm{LB}(R')\Rightarrow\mathrm{LB}(R)$.
\end{lemma}
\begin{proof}
For each $F$, apply the simulation to an optimal $R$-refutation, then take the
maximum over $F$. Polynomials are closed under composition.
\end{proof}

The bridge becomes \emph{easier} as $R$ gets stronger, and the lower bound
becomes \emph{harder}. A proof of $\mathrm P\ne\mathrm{NP}$ along this route
needs a single $R$ at which both hold.

\begin{proposition}[Reformulation trap]\label{prop:trap}
\status{PROVED/EXACT, given the cited KNOWN THEOREMS} If $\mathrm{LB}(R)$ is
already a theorem, then $\mathrm{Bridge}(R)$ is logically equivalent to
$\SAT\notin\mathrm P$. This applies to resolution (exponential lower bounds
for 3-CNF: Tseitin formulas on expanders \cite{urquhart1987}, random 3-CNF
\cite{chvatalszemeredi1988,bsw2001}), to tree-like resolution, to polynomial
calculus and to bounded-depth Frege on suitable 3-CNF families. The same
holds for any fixed restricted multisign-hypergraph calculus that is
p-simulated by one of these.
\end{proposition}
\begin{proof}
If the consequent of $\mathrm{Bridge}(R)$ is false, the implication is
equivalent to the negation of its antecedent.
\end{proof}

\begin{remark}
For such $R$, the expectation that ``the hard instances survive'' is
already settled and gives nothing new: the whole difficulty sits in the
bridge, and the bridge is exactly $\mathrm P\ne\mathrm{NP}$. A run that proves
a superpolynomial $\Wb_R$ for a resolution-strength hypergraph calculus has
made \emph{no} progress on the bridge. The world-ledger rule applies: a
statement equivalent to $\mathrm P\ne\mathrm{NP}$ is a reformulation, not
progress.
\end{remark}

\begin{proposition}[Trace-embedding: where free-sequence bounds do apply]\label{prop:trace}
\status{PROVED/EXACT} Let $A$ be a SAT algorithm. Suppose that from the run of $A$
on each $F\in\UNSAT$ one can read off, in time polynomial in the run, an
$R$-refutation of cost $\le q(\mathrm{time}_A(F))$. Then
$\mathrm{time}_A(F)\ge q^{-1}(w_R(F))$ for every such $F$. Examples: DPLL and
tree-like resolution; CDCL without restarts or with restarts and resolution
\cite{bks2004,pipatsrisawatdarwiche2011}.
\end{proposition}

This is the exact scope of the user's lower-bound strategy as it stands. It is
an unconditional lower bound for every algorithm whose \emph{own execution
contains} a derivation in $R$. An arbitrary polynomial-time $A$ need not do
this, and proving that it can always be forced to is $\mathrm{Bridge}(R)$ again.

\section{A bridge that is provable: reflection systems}

\begin{definition}[Reflection system $R_T$]
Let $T$ be a consistent, recursively axiomatized theory with polynomial-time
proof checking, for example $\mathrm{PA}$ or $\mathrm{ZFC}$, whose
$\Pi^0_1$ theorems are true. An $R_T$-refutation of $F$ is a tuple
$(\langle A\rangle,c,k,\rho,1^{m})$ with the following properties:
$\rho$ is a $T$-proof of
\[
 \Phi(A,c,k):\ \forall G\ \bigl[A(G)\text{ halts within }k(|G|+1)^c\text{ steps
 and outputs }\SAT(G)\bigr];
\]
$m=k(|F|+1)^c$; and $A$ rejects $F$ within $m$ steps. The extended system
$R_T^{+}$ also admits tuples $(\langle V\rangle,\rho,\sigma)$, where $\rho$ is a
$T$-proof that $V$ is a sound refutation verifier and $\sigma$ is a
$V$-refutation of $F$. To make both systems complete whatever $T$ proves,
each also admits every ordinary resolution refutation of $F$, verbatim.
\end{definition}

\begin{theorem}[Provable-correctness bridge]\label{thm:refl}
\status{PROVED/EXACT}
\begin{enumerate}
\item $R_T$ and $R_T^{+}$ are sound and complete Cook--Reckhow systems. Completeness comes from the resolution clause. $\Phi$ is
$\Pi^0_1$, so it is true when provable. The verifier checks $\rho$ and then
simulates $A$ for $m$ steps; both take time polynomial in the tuple length,
and $1^m$ is part of the input.
\item \emph{(Bridge.)} If some $A$ with $T\vdash\Phi(A,c,k)$ exists, then
$\Wb_{R_T}(N)\le |A|+|\rho|+O(\log k+\log c)+k(N+1)^c=O(N^c)$. Here
$|A|$ and $|\rho|$ are constants independent of $N$.
\item \emph{(Strength.)} $R_T^{+}$ p-simulates every proof system whose
soundness $T$ proves, with an additive constant overhead. For
$T\supseteq S^1_2$ this includes extended Frege \cite{cook1975,krajicekpudlak1989}.
\item Hence, by Lemma~\ref{lem:mono}, $\mathrm{LB}(R_T^{+})\Rightarrow\mathrm{LB}(EF)$,
and indeed $\mathrm{LB}$ for every $T$-provably sound system.
\end{enumerate}
\end{theorem}

\begin{corollary}[Exact location of the gap]\label{cor:gap}
\status{PROVED/EXACT as an implication; both premises OPEN}
\[
 \mathrm{LB}(R_T)\ \Longrightarrow\ \text{no polynomial-time SAT algorithm is $T$-provably correct and $T$-provably polynomial.}
\]
Upgrading this to $\SAT\notin\mathrm P$ needs exactly one extra statement:
\[
 \mathrm{PROV}_T:\ \SAT\in\mathrm P\ \Longrightarrow\ \exists A,c,k:\ T\vdash\Phi(A,c,k).
\]
\end{corollary}

So the weakest \emph{sufficient} bridge splits into a part that is free
(Theorem~\ref{thm:refl}.2) and an independence question (PROV$_T$). PROV$_T$ is
not known, and its failure is consistent with present knowledge: P could equal
NP via an algorithm whose correctness $T$ cannot prove. Meanwhile
$\mathrm{LB}(R_T)$ is at least as hard as superpolynomial extended Frege lower
bounds, which are a central open problem. The squeeze is therefore real.
Strengthening $R$ until the bridge becomes provable moves the whole
difficulty into the lower bound.

\section{Literature check this run}

\begin{table}[h]
\centering\small
\begin{tabular}{p{4.4cm}p{3.1cm}p{7.2cm}}
\toprule
Source & Class & Exact relevance\\
\midrule
Cook--Reckhow 1979 \cite{cookreckhow1979} & KNOWN THEOREM & A p-bounded system exists iff NP $=$ coNP. Existential: it gives no bridge for a fixed $R$.\\
Kraj\'{\i}\v{c}ek--Pudl\'ak 1989 \cite{krajicekpudlak1989} & KNOWN THEOREM & E $=$ NE $\Rightarrow$ a p-optimal system exists; NE $=$ coNE $\Rightarrow$ an optimal system exists. P $=$ NP $\Rightarrow$ E $=$ NE, so under P $=$ NP \emph{some} system is p-bounded \emph{and} p-optimal, but not an explicitly named $R$. Our $R_T$ is the explicit, provability-relativized form.\\
K\"obler--Messner--Tor\'an 2003 & KNOWN THEOREM & Weakens these hypotheses to EE $=$ NEE and NEE $=$ coNEE.\\
Pich--Santhanam 2023 \cite{pichsanthanam2023} & KNOWN THEOREM (conditional) & The closest known result to Bridge(EF). Under two $S^1_2$-provability hypotheses (hardness of a concrete E function, and a reduction from breaking one-way functions to learning), a circuit-lower-bound family that is hard for EF yields P $\ne$ NP. Several further hypotheses each give ``EF not p-bounded $\Rightarrow$ P $\ne$ NP''. This is exactly the shape of Bridge(EF), conditional on provability. It matches our Corollary~\ref{cor:gap} in kind: provability in bounded arithmetic is the missing link.\\
Pich--Santhanam FOCS 2019 \cite{pichsanthanam2019} & KNOWN THEOREM (barrier) & An analogue of natural proofs for proof complexity. Under hardness assumptions, strong systems cannot be shown to have lower bounds for natural random distributions of hard tautologies. This is a hostile control for the Hypergraph \& Certificates role.\\
Buss--Williams 2015 \cite{busswilliams2015} & KNOWN THEOREM (barrier) & World-ledger row L13. It limits the time--space spine, not $\Wb_R$.\\
\bottomrule
\end{tabular}
\end{table}

\paragraph{Barrier-transfer classification.}
(i) Pich--Santhanam: a \emph{directly transferable theorem} for $R=EF$
(conditional). (ii) Kraj\'{\i}\v{c}ek--Pudl\'ak optimality: a \emph{structural analogy}.
It needs an \emph{explicit} p-optimal $R$, which is OPEN. (iii) Post-quantum
technique \#5 (worst-case/average-case extraction) corresponds to
``extract an $R$-proof from black-box solver behaviour''. That is exactly
PROV$_T$ or a trace embedding, so it is an ANALOGY/CANDIDATE with the same
gap. (iv) Control theory (barrier certificates): the dual object to a short
$R$-refutation is a lower-bound certificate, such as a width, degree or
Prover--Delayer strategy. It transfers exactly only for systems where a
game characterization exists (resolution, polynomial calculus), so the transfer
fails for strong $R$.

\section{Red-team checks on the shared target}
\begin{enumerate}
\item \emph{Cost/measure mismatch.} If $\cost_R$ counts only DAG nodes, but
nodes can carry exponentially large signed hyperedges, then $\Wb_R$ can be
small while the encoding is large. Rule: include encoded size, or prove a
polynomial relation between the two. \status{EXACT requirement}
\item \emph{Space variant.} If peak memory is the measure, the bridge must be
SAT $\in$ some space class $\Rightarrow$ small space measure. The Prop.~\ref{prop:trap}
trap applies again: resolution clause-space lower bounds are known
\cite{bsg2003}.
\item \emph{Adaptive strategy switching.} A fixed $R$ that contains several
rule families, such as resolution, algebraic and circuit-compression rules, is
p-simulated by their combination. Its lower bound needs a technique that
works against \emph{all} rule families at once. No such technique is known
above bounded-depth Frege for 3-CNF. \status{OPEN}
\item \emph{Upper bounds help the red team.} Every explicit short derivation,
such as an LSPV W61 certificate, lowers $w_R$. A claimed lower bound
must be checked against them.
\end{enumerate}

\section{WHAT CHANGED RELATIVE TO ROOT MAIN}
Nothing in root \texttt{main.tex} is contradicted. This run adds the following
points, for LSPV-5 to absorb if it chooses. (a) Any $\Wb_R$ lower bound for a
resolution-strength calculus is \emph{not} progress on the bridge, by
Prop.~\ref{prop:trap}. (b) There is an explicit system with a provable
\emph{conditional} bridge (Theorem~\ref{thm:refl}). (c) The single residual
statement for that system is PROV$_T$.

\section{ADOPTED / REJECTED / RETIRED PRIOR COMPONENTS}
Adopted: the Codex 12:03 $W_R$ definition and its insistence on fixed $A,c$.
Rejected as a bridge, though kept as true: ``Cook--Reckhow gives the bridge''.
It is existential. Not retired: the user's free-sequence lower-bound principle,
which is valid and scoped by Prop.~\ref{prop:trace}.

\section{STALE OR OVERSTATED CLAIMS FOUND}
None in root \texttt{main.tex}. The scheduled prompt cites a ``shared worst-case
research directive on GitHub''. No such file existed before this run, and it is
now created. That was a stale reference, not an overstatement.

\section{Five-role assignments for the next runs}
\begin{enumerate}
\item \textbf{General SAT Lower-Bound Lead.} Prove or refute PROV$_T$ for
$T=\mathrm{ZFC}$ in a relativized world. The first step is to construct an
oracle $O$ with $\mathrm P^O=\mathrm{NP}^O$ but with no $T$-provably correct
polynomial-time SAT$^O$ algorithm. Such an oracle would show that
Corollary~\ref{cor:gap} cannot be closed by relativizing methods.
\item \textbf{SAT-Circuit Reduction Builder.} Fix the combined multisign
hypergraph calculus $R_{\mathrm{mix}}$ explicitly, with rules, cost and
encoding. Then place it in the simulation order, either at most bounded-depth
Frege (trap applies) or at least Frege (bridge question nontrivial).
\item \textbf{Hypergraph \& Proof Certificates.} Use Pich--Santhanam 2019 as
the hostile control. For each hard UNSAT family, record whether it is a
``natural'' random distribution, which is barrier-exposed.
\item \textbf{Time--Space Resource Auditor.} Check that $\cost_R$ is
polynomially related to encoded proof length for $R_{\mathrm{mix}}$. Verify the
$R_T$ verifier runtime claim in Theorem~\ref{thm:refl}.1.
\item \textbf{Independent Proof Red Team.} For each claimed $\Wb_R$ lower
bound, ask whether the trap (Prop.~\ref{prop:trap}) applies. For each claimed
bridge, ask whether it silently assumes PROV$_T$ or a trace embedding.
\end{enumerate}

\section{Status table}
\begin{center}\small
\begin{tabular}{lll}
\toprule
Statement & Status & Note\\
\midrule
Free sequence choice is sound for lower bounds & PROVED/EXACT & minimum over sequences\\
Lemma~\ref{lem:mono} monotone squeeze & PROVED/EXACT & folklore\\
Prop.~\ref{prop:trap} reformulation trap & PROVED/EXACT & trivial logic plus known LBs\\
Prop.~\ref{prop:trace} trace embedding & PROVED/EXACT & folklore; CDCL/resolution\\
Theorem~\ref{thm:refl} reflection bridge & PROVED/EXACT & folklore-type (KP89)\\
PROV$_T$ & OPEN & independence-type\\
$\mathrm{LB}(R_T)$, $\mathrm{LB}(EF)$ & OPEN & central open problem\\
Bridge$(R)$ for any $R$ with known LB & equivalent to P $\ne$ NP & reformulation\\
P versus NP & OPEN & \\
\bottomrule
\end{tabular}
\end{center}

\begin{thebibliography}{99}
\bibitem{cookreckhow1979} S.~A.~Cook and R.~A.~Reckhow, The relative efficiency of propositional proof systems, \emph{J.\ Symbolic Logic} 44 (1979) 36--50.
\bibitem{cook1975} S.~A.~Cook, Feasibly constructive proofs and the propositional calculus, STOC 1975.
\bibitem{krajicekpudlak1989} J.~Kraj\'{\i}\v{c}ek and P.~Pudl\'ak, Propositional proof systems, the consistency of first order theories and the complexity of computations, \emph{J.\ Symbolic Logic} 54 (1989) 1063--1079.
\bibitem{pichsanthanam2023} J.~Pich and R.~Santhanam, Towards P $\ne$ NP from Extended Frege lower bounds, arXiv:2312.08163; ECCC TR23-199. \url{https://arxiv.org/abs/2312.08163}
\bibitem{pichsanthanam2019} J.~Pich and R.~Santhanam, Why are proof complexity lower bounds hard?, FOCS 2019.
\bibitem{busswilliams2015} S.~Buss and R.~Williams, Limits on alternation trading proofs for time-space lower bounds, \emph{Comput.\ Complexity} 24 (2015).
\bibitem{urquhart1987} A.~Urquhart, Hard examples for resolution, \emph{J.\ ACM} 34 (1987).
\bibitem{chvatalszemeredi1988} V.~Chv\'atal and E.~Szemer\'edi, Many hard examples for resolution, \emph{J.\ ACM} 35 (1988).
\bibitem{bsw2001} E.~Ben-Sasson and A.~Wigderson, Short proofs are narrow---resolution made simple, \emph{J.\ ACM} 48 (2001).
\bibitem{bsg2003} E.~Ben-Sasson and N.~Galesi, Space complexity of random formulae in resolution, \emph{Random Struct.\ Alg.} 23 (2003).
\bibitem{bks2004} P.~Beame, H.~Kautz and A.~Sabharwal, Towards understanding and harnessing the potential of clause learning, \emph{JAIR} 22 (2004).
\bibitem{pipatsrisawatdarwiche2011} K.~Pipatsrisawat and A.~Darwiche, On the power of clause-learning SAT solvers as resolution engines, \emph{Artif.\ Intell.} 175 (2011).
\end{thebibliography}

\end{document}
% END_OF_VERIFIED_REPORT::CLAUDE::WORSTCASE-BRIDGE-0001
